% ============================================================
\chapter{Introduction}
% ============================================================

\section{Purpose of the Project}

Modern computer networks combine a wide range of devices, including routers,
switches, firewalls, and Internet of Things equipment such as IP cameras.
Administering and monitoring these devices normally requires operators to be
familiar with several vendor-specific command-line interfaces, which makes
network operations slow, error-prone, and difficult for less experienced
staff to perform safely.

The purpose of this project is to design and implement NetMind, a network
operations assistant that allows a user to manage and monitor a network using
natural language instead of memorised commands. NetMind is built around a
Large Language Model, accessed through the Groq API, which is given a set of
callable tools that can check a device's power state and IP address, turn
devices on or off, list every device on a given network segment, test whether
two devices can reach each other, audit a firewall's configuration for common
security weaknesses such as an open SSH port, and verify whether a security
camera is still using its default password. The system understands requests
in both Arabic and English and responds through a web application that
combines the chat with a live monitoring dashboard.

To allow safe experimentation, NetMind operates in a simulated mode against a
virtual network built with Containerlab, consisting of a router, firewalls,
switch and hub segments, hosts, and cameras running inside Docker containers
with real IP addressing and routing between branches. A second, read-only
real-network mode extends the same natural-language interface to an actual
physical network, discovering connected devices by combining ARP scanning, a
ping sweep, IPv6 neighbour discovery and mDNS. Together, these two modes
demonstrate that a single conversational interface can be used both to learn
and test network administration safely, and to observe a real network
without requiring the user to run separate diagnostic tools.

\section{Purpose of this Document}

This document is prepared in partial fulfilment of the requirements for the
degree of Bachelor of Science in Computer Science at the College of Computer
Engineering and Sciences, Prince Sattam Bin Abdulaziz University. It records
the motivation behind NetMind, the analysis and design decisions taken while
building it, the implementation of its simulated and real-network modes, the
testing carried out to confirm that the system behaves correctly and safely,
and the plan for the second term of the project.

The primary readers of this document are the project supervisor,
Dr.\ Khalil Chebil, and the examining committee responsible for assessing
this Final Year Project. It is also intended to be useful to any student or
developer who wishes to understand, maintain, or extend NetMind in the
future, by explaining not only what was built but why particular design
choices were made.

\section{Overview of this Document}

The remainder of this document is organised as follows.

Chapter~2, Literature Review, examines existing approaches to network
management and related tools, and positions NetMind relative to them.

Chapter~3, System Analysis, presents the requirements gathered for the
project, including functional and non-functional requirements, together with
a diagram describing how a user interacts with the system.

Chapter~4, Design Considerations, discusses the constraints that shaped the
project, the hardware and software environment used, and the development
method followed.

Chapter~5, System Design, describes the architecture of NetMind in detail:
its layered structure, the device registry and services, the tool registry
that connects them to the Large Language Model, the simulated Containerlab
network, the real-network discovery mode and the live monitoring dashboard,
together with its confirmation and audit safeguards, its persistent storage,
and the testing carried out.

Chapter~6, Conclusion, summarises what was achieved, reflects on the
limitations of the current system, and sets out the work planned for the
second term, followed by the References and Appendices.
